\scanpage{22}\mspage{- 17 -}

et la loi de composition par
\[
(u',s',t')\mathbin{\unsure{o|\circ}}(u,s,t)=(u\,u',\,u's+t',\,u^{\prime 3}t+t'+u^{\prime 2}rs')\ .
\]

Lorsque la courbe est mise sous forme (6.8), ces équations se réduisent à
\[
\begin{aligned}
u\in\mu_3\qquad s\in\mathbb{F}_4\qquad t^2+t=s^3\quad &(=0\text{ ou }1\text{ selon que}\\
&s=0\text{ ou }s\ne0)\ .
\end{aligned}
\]

7. \uline{Anneau des formes modulaires entières}.

On appelle forme modulaire entière de poids $n$ une loi qui à chaque courbe de genre un sur une base $S$ associe une section de $\omega^{\otimes n}$, et ce de façon compatible au changement de base. Il reviendrait au même de se limiter aux courbes sans fibres additives.

\uline{Proposition} 7.1. \uline{L'anneau des formes modulaires entières est engendré sur} $\mathbb{Z}$ \uline{par} $c_4$, $c_6$ \uline{et} $\Delta$ \uline{soumis à la seule relation}
\[
c_4^3-c_6^2=1728\,\Delta\quad.\tag*{(7.2.)}
\]

Appliquant la définition à la courbe (1.3), on voit qu'une forme modulaire \unsure{\.{e}ntière|éntière|entière} de poids $n$ s'exprime comme polynôme de\margin{\handwritten{$\unsure{\sigma|6|o}$ ?}} poids $n$ en les $a_i$. Pour $2$ et $3$ inversibles, il résulte de 2.5 que ces formes sont les polynômes en $c_4$ et $c_6$, et en particulier que $c_4$, $c_6$ et $\Delta$ ne satisfont pas d'autres relations que (7.2). Reste à prouver que si un polynôme en $c_4$, $c_6$ et $\Delta$ est identiquement nul en caractéristique $2$ (resp. $3$), alors il est divisible par $2$ (resp.$3$) dans $\mathbb{Z}[c_4,c_6,\Delta]/c_4^3-c_6^2-1728\,\Delta$, ce qui signifie

% archive sheet number: 21
